\documentclass[runningheads,envcountsame,envcountsect]{llncs}

\let\vec\relax
\usepackage{amsmath,amssymb,mathtools}
\usepackage{booktabs}
\usepackage{microtype}
\usepackage[hidelinks]{hyperref}
\usepackage{cite}
\hypersetup{
  pdftitle={Changing Characteristic in Fast Linear Codes for Pseudorandom
    Correlation Generators},
  pdfauthor={Stanislav Peceny and Peter Rindal}
}

\newcommand{\F}{\mathbb F}
\newcommand{\E}{\mathbb E}
\newcommand{\wt}{\operatorname{wt}}
\newcommand{\supp}{\operatorname{supp}}
\newcommand{\rank}{\operatorname{rank}}
\newcommand{\cC}{\mathcal C}
\newcommand{\cF}{\mathcal F}
\newcommand{\cZ}{\mathcal Z}
\newcommand{\negl}{\operatorname{negl}}

\title{Changing Characteristic in Fast Linear Codes for Pseudorandom
Correlation Generators}
\titlerunning{Changing Characteristic in Fast Linear Codes}
\author{Stanislav Peceny \and Peter Rindal}
\authorrunning{S. Peceny and P. Rindal}
\institute{}

\begin{document}
\raggedbottom
\maketitle

\begin{abstract}
Several pseudorandom correlation generators use linear codes whose encoders
were designed over the binary field.  Applying the same addition circuit to
vectors over a large prime field avoids general field multiplications, but the
binary distance analysis need not survive the change of characteristic.  This
paper investigates when such a transfer is justified.  Scalar extension is
benign: a matrix over a subfield has the same minimum distance over every
extension field.  Changing from characteristic two to odd characteristic is
different, and explicit addition-only encoders have smaller distance over an
odd prime field even when the prime exceeds the block length.  The failure is
not an impossibility result.  A dense random binary generator has constant
rate and relative distance close to one half over the binary field and any
prescribed finite family of prime fields.  The construction uses only
additions, but direct encoding is quadratic.

The main case study is the rate-$1/3$ repetition-three BAA ensemble.  We prove
relative-distance lower bounds $0.196$, $0.05$, and $0.018$ over
$\F_3$, $\F_5$, and $\F_7$, respectively.  More generally, the ensemble has
positive relative distance over every fixed prime field.  For a prime that
grows with the security parameter, we prove an absolute distance floor and a
cycle-rank barrier for characteristic-specific dependencies; a linear-distance
theorem remains open.  We also examine Walsh--Hadamard codes.  They provide a
direct large-field approach, but restricting every spectral multiplier to a
fixed alphabet forces vanishing relative distance.  These results separate
existence from efficiency and identify a concrete open problem: construct a
near-linear, addition-only encoder with a field-uniform linear-distance proof.
\keywords{linear codes \and pseudorandom correlation generators \and
learning parity with noise \and minimum distance \and prime fields}
\end{abstract}

\section{Introduction}\label{sec:introduction}

Pseudorandom correlation generators (PCGs) reduce communication by replacing
a long correlated string with short seeds and local expansion.  The expansion
often contains a multiplication by a public generator matrix.  This operation
is inexpensive over $\F_2$, where every nonzero coefficient is one.  Over a
large prime field, a general matrix requires many field multiplications.

A tempting optimization is to retain a binary generator matrix and reinterpret
its entries in the target field.  The resulting encoder computes only sums of
selected message coordinates.  This optimization raises a basic question:

\begin{quote}
When does a binary matrix retain its minimum-distance guarantee after the
message and codeword alphabets change to $\F_p$?
\end{quote}

The question has no unconditional affirmative answer.  Reduction modulo a new
prime changes the ranks of integer minors and creates new cancellation
patterns.  Conversely, large alphabets can remove binary cancellations and
increase distance.  Neither direction is monotone for a fixed encoder.

The literature contains several different responses to this problem.
Unique-neighbor expansion gives field-independent distance for some sparse
matrices~\cite{silver21}.  Analyses of expand-accumulate and
expand-convolute codes are instead tied to binary weight processes
\cite{ea22,ec23}.  Security analyses may compensate for a smaller distance by
increasing the noise weight, but such compensation does not prove a code
theorem~\cite{silver21,stationary25}.  Recent Walsh--Hadamard PCGs avoid a
binary transfer altogether by analyzing a code directly over the deployed
prime field~\cite{wht26}.

This paper organizes these mechanisms around the change-of-characteristic
problem.  The objective is not to defend a universal inheritance heuristic.
It is to determine what can be proved, what fails, and what an efficient
construction must supply.

\paragraph{Contributions.}
The paper gives the following results.

\begin{enumerate}
  \item It separates scalar extension from change of characteristic.  Scalar
  extension preserves minimum distance exactly; change of characteristic does
  not.

  \item It gives a positive baseline.  A random dense binary generator has
  constant rate and relative distance close to $1/2$ simultaneously over any
  prescribed finite family of fields.  The failure bound is independent of
  the field sizes.

  \item It records explicit negative examples.  Samplewise distance equality
  fails for addition-only BAA encoders, including an example with $p>N$.
  A separate family shows that unrestricted binary matrices can have binary
  relative distance bounded away from zero and vanishing ternary relative
  distance, although that family has vanishing rate.

  \item It proves that repetition-three BAA codes are asymptotically good over
  every fixed prime field.  The ternary bound $0.196$ exceeds the published
  binary lower bound $0.1323$ for the same two-accumulator ensemble
  \cite{kliewer08}.

  \item It proves two barriers to failures over large fields.  In BAA, a
  characteristic-$p$ distance loss requires a logarithmic-size support and a
  residual component with at least $\lceil\log_2p\rceil$ independent cycles.

  \item It shows that a Walsh--Hadamard encoder with a fixed random spectral
  alphabet has vanishing relative distance.  A growing alphabet remains a
  possible finite-parameter design.
\end{enumerate}

The results locate the remaining difficulty.  Good binary matrices over prime
fields exist, and fast addition-only encoders exist.  No present construction
combines a near-linear encoder with a field-uniform distance proof at useful
rate and distance.

\paragraph{Organization.}
Section~\ref{sec:fields} distinguishes field changes and reviews the security
role of distance.  Section~\ref{sec:baselines} gives the dense positive
baseline and the negative examples.  Section~\ref{sec:baa} develops the BAA
case study.  Section~\ref{sec:wht} examines the Walsh--Hadamard alternative.
Section~\ref{sec:discussion} states the resulting design requirements and open
problems.

\section{Field changes and their security role}\label{sec:fields}

Let $F$ be a field and let $G\in F^{k\times n}$.  The code generated by $G$ is

\[
  \cC_F(G):=\{xG:x\in F^k\}.
\]

Its minimum distance is

\[
  d_F(G):=\min\{\wt(xG):x\in F^k,\ xG\ne0\}.
\]

The matrix need not have full row rank for this definition.  The code rate is
$\rank_F(G)/n$.

\subsection{Scalar extension is exact}

Suppose $K/F$ is a field extension and $G$ has entries in $F$.  The same
matrix defines codes over $F$ and $K$.

\begin{theorem}[Scalar-extension invariance]\label{thm:extension}
For every field extension $K/F$ and every matrix $G\in F^{k\times n}$,
\[
  d_K(G)=d_F(G).
\]
\end{theorem}

\begin{proof}
Fix an $F$-basis $\beta_1,\ldots,\beta_s$ of the finite-dimensional subspace
of $K$ spanned by the coordinates of a message $x\in K^k$.  Write
$x=\sum_i\beta_i x_i$ with $x_i\in F^k$.  Then
\[
  xG=\sum_i\beta_i(x_iG).
\]
A coordinate of $xG$ is zero exactly when the same coordinate is zero in every
$x_iG$.  Hence the support of $xG$ is the union of the supports of its nonzero
$F$-components.  Every nonzero $K$-codeword therefore has weight at least
$d_F(G)$.  The reverse inequality follows from $\cC_F(G)\subseteq\cC_K(G)$.
\end{proof}

Theorem~\ref{thm:extension} covers, for example, a binary matrix interpreted
over $\F_{2^s}$.  It does not cover $\F_p$ for odd $p$, because no field
embedding $\F_2\hookrightarrow\F_p$ exists.

\subsection{Changing characteristic}

Let $G_{\mathbb Z}\in\mathbb Z^{k\times n}$ be an integer matrix.  Reduction
modulo a prime $p$ gives a matrix $G_p$.  A selected column set can be
independent over $\F_2$ and dependent over $\F_p$ when its nonzero integer
minors are divisible by $p$.  Thus the obstruction is torsion in an integer
cokernel, not a failure of scalar extension.

The same distinction appears at the codeword level.  A binary sum vanishes
when an even number of active terms meet.  Over $\F_p$, the values of those
terms matter.  New cancellations can appear, while old parity cancellations
can disappear.

\subsection{Distance, noise, and PCG security}

Minimum distance is not itself an LPN assumption.  It prevents linear tests
from finding an unusually short dual relation and supports parameter choices
against known attacks.  Silver makes both points explicit.  Its
unique-neighbor argument applies to binary sparse matrices over arbitrary
fields, while its security discussion observes that additional noise can
compensate for an overestimate of distance against linear attacks
\cite{silver21}.

Expand-Convolute demonstrates why empirical extrapolation is insufficient.  It
gives efficient attacks against the proposed Silver codes and develops a new
binary code with a direct distance proof~\cite{ec23}.  Its Markov-chain
analysis is binary; it is not a change-of-characteristic theorem.

Stationary Syndrome Decoding gives a complementary security result.  Against
linear tests, its bias depends explicitly on code distance and noise weight,
and a larger noise weight can offset a smaller distance in that model
\cite{stationary25}.  This tradeoff does not establish the distance of the
matrix over a new field.  A code theorem and a security compensation argument
answer different questions.

\section{Positive and negative baselines}\label{sec:baselines}

The field-transfer problem is meaningful only if some binary matrices work
over many characteristics.  This section first gives such a construction.  It
then shows why binary distance alone does not certify an arbitrary matrix.

\subsection{A dense field-universal code}

Let $G_{\square}$ be the $k$-by-$2^k$ binary matrix whose columns are all
vectors in $\{0,1\}^k$.

\begin{proposition}[Boolean evaluation code]\label{prop:cube}
For every field $F$, the code $\cC_F(G_{\square})$ has dimension $k$ and
minimum distance $2^{k-1}$.
\end{proposition}

\begin{proof}
Fix $x\in F^k\setminus\{0\}$ and an index $i$ with $x_i\ne0$.  Pair each
$a\in\{0,1\}^k$ with the vector obtained by flipping coordinate $i$.  Their
inner products with $x$ differ by $x_i$, so at most one inner product in each
pair is zero.  Every nonzero word has weight at least $2^{k-1}$.  A message
supported on one coordinate attains this weight.
\end{proof}

The rate of this code vanishes.  Sampling only linearly many columns recovers
constant rate.

For a field $F$ and nonzero $x\in F^k$, define
\[
  Z_x:=\{a\in\{0,1\}^k:\langle x,a\rangle=0\}.
\]
The proof of Proposition~\ref{prop:cube} gives $|Z_x|/2^k\le1/2$.

\begin{lemma}[Kernel sections]\label{lem:vc}
The set system $\cZ_F:=\{Z_x:x\in F^k\setminus\{0\}\}$ has VC dimension at
most $k-1$.
\end{lemma}

\begin{proof}
Suppose $a_1,\ldots,a_m$ are shattered.  One labeling places every $a_j$ in
the kernel of a nonzero functional, so their span is proper.  For each $j$,
another labeling contains every $a_i$ with $i\ne j$ but excludes $a_j$.
Hence no $a_j$ lies in the span of the other points.  The points are linearly
independent and span a proper subspace, which implies $m\le k-1$.
\end{proof}

Sample independent columns $A_1,\ldots,A_n\gets\{0,1\}^k$ and set
$G=[A_1|\cdots|A_n]$.  A standard VC inequality and
Lemma~\ref{lem:vc} give the following theorem~\cite{vc71}.

\begin{theorem}[Dense simultaneous baseline]\label{thm:dense}
Fix fields $F_1,\ldots,F_t$, a real $0<\varepsilon<1/2$, and integers
$k\ge2$ and $n$ with $2n\ge k-1$.  Except with probability at most
\[
  8t\left(\frac{2en}{k-1}\right)^{k-1}
      \exp\left(-\frac{n\varepsilon^2}{32}\right),
\]
the sampled matrix $G$ satisfies, simultaneously for every $s\in[t]$,
\[
  \rank_{F_s}(G)=k
  \quad\text{and}\quad
  d_{F_s}(G)\ge(1/2-\varepsilon)n.
\]
\end{theorem}

\begin{proof}
Apply the VC inequality to the uniform distribution on $\{0,1\}^k$ and the
class $\cZ_{F_s}$.  Uniformly for every nonzero $x\in F_s^k$, the empirical
fraction of columns in $Z_x$ is at most $1/2+\varepsilon$ outside the stated
failure event.  The word $xG$ therefore has weight at least
$(1/2-\varepsilon)n$.  Its positive weight also proves full row rank.  A union
bound over the $t$ fields completes the proof.
\end{proof}

For fixed $\varepsilon$, a sufficiently small constant rate makes the failure
probability exponentially small in $k$.  Theorem~\ref{thm:dense} also handles
growing primes.  Taking all primes below $2^\lambda$ adds at most $\lambda$
bits to the union bound.  If $k=\Omega(\lambda)$, one binary matrix can retain
constant rate and distance over $\F_2$ and every such prime field.

Direct encoding computes $n$ dense subset sums and costs $\Theta(kn)$ field
additions.  At constant rate this is quadratic.  Theorem~\ref{thm:dense}
settles existence but not efficient encoding.

\subsection{Why binary distance is insufficient}

An explicit family gives a minimal warning.  Let $M:=J_4-I_4$.  For $L\ge1$,
form a $4$-by-$(4L+1)$ binary matrix $G_L$ by repeating each column of $M$
$L$ times and appending one all-one column.  Since $\det(M)=-3$,
\[
  d_{\F_2}(G_L)=L+1,
  \qquad
  d_{\F_3}(G_L)=1.
\]
Both reductions have row rank four.  Consequently,
\[
  \frac{d_{\F_2}(G_L)}{4L+1}\longrightarrow\frac14,
  \qquad
  \frac{d_{\F_3}(G_L)}{4L+1}\longrightarrow0.
\]
The family has vanishing rate, so it does not rule out a theorem with an
explicit positive-rate hypothesis.  It does rule out an inheritance theorem
for unrestricted binary matrices based only on their binary relative distance.

The BAA case gives positive-rate finite counterexamples.  For length six, one
rate-$1/3$ repetition-three encoder has $d_{\F_2}=4$ and $d_{\F_7}=3$.
Thus even $p>N$ does not guarantee samplewise monotonicity.  These examples
motivate ensemble statements rather than equality for every realized encoder.

\section{Addition-only BAA codes}\label{sec:baa}

Block-accumulate-accumulate codes were introduced as fast linear encoders for
PCGs and zero-knowledge systems~\cite{baa26}.  This section studies their
addition-only repetition specialization.

Fix $N=3K$.  Let $B_3:F^K\to F^N$ repeat each message coordinate three
times, let $A$ be the prefix-sum accumulator, and sample independent uniform
coordinate permutations $P_1,P_2$.  The encoder is
\[
  E_{P_1,P_2}:=A P_2 A P_1 B_3.
\]
It has rate $1/3$ over every field and uses only additions and data movement.
The same permutations define a coupled code over every characteristic.

The first-moment analysis tracks full symbol compositions through both
accumulators.  A separate exact spacing argument controls messages whose
support is too small for normalized entropy estimates.  Together they give
the following fixed-prime results.

\begin{theorem}[Fixed-prime BAA bounds]\label{thm:baa-fixed}
For the repetition-three BAA ensemble,
\begin{align*}
  \Pr[d_{\F_3}(E_{P_1,P_2})\le \lfloor 0.196N\rfloor]&\longrightarrow0,\\
  \Pr[d_{\F_5}(E_{P_1,P_2})\le \lfloor N/20\rfloor]&\longrightarrow0,\\
  \Pr[d_{\F_7}(E_{P_1,P_2})\le \lfloor 9N/500\rfloor]&\longrightarrow0.
\end{align*}
The probability is over $(P_1,P_2)$ as $K\to\infty$.
\end{theorem}

The ternary bound is stronger than the binary lower bound $0.1323$ reported
for repetition three and two accumulators~\cite{kliewer08}.  It therefore
proves one-sided threshold inheritance throughout the certified binary range.
It does not prove equality of the binary and ternary limiting distances.

The exact enumerator also suggests a sharper prediction.  For a fixed prime
$p$, let $\Phi_p(\tau)$ denote the exponential rate contributed by an
admissible vector-composition and transition type $\tau$.  Let $a(\tau)$ and
$c(\tau)$ be its normalized message and output weights.  Define the
first-moment threshold
\[
  \delta_{\rm FM}(p):=\inf D_p,
  \qquad
  D_p:=\left\{\delta>0:
  \begin{array}{l}
  \text{some admissible $\tau$ satisfies $a(\tau)>0$,}\\[-1mm]
  \text{$c(\tau)\le\delta$, and $\Phi_p(\tau)>0$}
  \end{array}\right\}.
\]
The restriction $a(\tau)>0$ isolates linear-weight messages.  The spacing
argument treats message weights that vanish after normalization.

Numerical optimization gives the following sign brackets.  These values are
computational evidence, not certified global optima.
\begin{center}
\begin{tabular}{@{}ccc@{}}
\toprule
field & negative exponent & positive exponent \\
\midrule
$\F_2$ & $\delta=0.132$ & $\delta=0.133$ \\
$\F_3$ & $\delta=0.195$ & $\delta=0.197$ \\
$\F_5$ & $\delta=0.260600$ & $\delta=0.260612$ \\
\bottomrule
\end{tabular}
\end{center}
The binary bracket reproduces the known numerical transition.  The odd-prime
brackets predict a strict improvement rather than equality with the binary
distance.

For $p=5$, the full transition search can be reduced without relaxing its
objective.  Let $\alpha$ be an increment composition and let $\beta$ be a
stationary state composition.  For a law $R$ on $(D,S)\in\F_5^2$, require

\[
 D\sim\alpha,\qquad S\sim\beta,\qquad D+S\sim\beta.
\]

Set $U:=\mathbf 1[D\ne0]$ and $V:=\mathbf 1[S\ne0]$.  Fix the crossing
$x:=\Pr[U=1,V=0]$, and let $J_x$ be the resulting law of $(U,V)$.  The
conditional symbol laws are

\[
 \alpha_u(d):=\Pr_{\alpha}[D=d\mid U=u],
 \qquad
 \beta_v(s):=\Pr_{\beta}[S=s\mid V=v].
\]

The support-preserving product law is

\[
 Q_x(d,s):=J_x(u,v)\alpha_u(d)\beta_v(s),
 \quad u=\mathbf 1[d\ne0],\quad v=\mathbf 1[s\ne0].
\]

For every feasible $R$ with crossing $x$, the chain rule for relative entropy
gives

\[
 D_{\rm KL}(R\,\|\,\alpha\otimes\beta)
 =I_{J_x}(U;V)+D_{\rm KL}(R\,\|\,Q_x).
\]

Thus the exact stationary information payment equals

\[
 I_{\rm stat}(\alpha,\beta)
 =\min_x\left(I_{J_x}(U;V)
   +\min_R D_{\rm KL}(R\,\|\,Q_x)\right),
\]

where the inner minimum uses the three displayed marginals and the fixed
support table.  Its entropy dual has thirteen potentials at interior quinary
points.  Applying the reduction to both accumulators replaces two $5$-by-$5$
transition laws with three compositions and two crossings.

As a numerical fidelity check, the reduced objective and the full transition
objective agree within $1.9\cdot10^{-7}$ at output fractions $0.22$, $0.24$,
$0.25$, $0.26$, $0.260600$, $0.260612$, and $0.261$.  The discrepancy is
below $1.2\cdot10^{-13}$ at the last four points.  These calculations explain
the refined quinary sign bracket in the table.  They do not certify that the
heuristic searches found the global maximum.

The stationary payment also admits a useful fixed-dual lower bound directly
in the two composition laws.  A feasible entropy dual determines constants
$c,u$ and a positive matrix $K$ such that

\[
 I_{\rm stat}(\alpha,\beta)
 \ge c+\langle u,(\alpha,\beta)\rangle
       -\log(\alpha^{\mathsf T}K\beta).
\]

The dual data are fixed before the compositions are optimized.  A finite
bank of such inequalities can therefore remain active at the same middle
composition in both accumulator terms.  This avoids the quantifier loss that
arises if the two payments minimize over the middle composition separately.

We used this envelope to audit two points below the predicted quinary
transition.  At output fraction $0.22$, a bank of 72 cuts per accumulator
produced no positive local maximum on a saved partition of 980 composition
boxes.  At output fraction $0.24$, a separate bank of 93 cuts produced no
positive local maximum on a 101-box partition; the largest value found was
$-5.93\cdot10^{-4}$.  The latter scan records two successful optimization
starts on each box not already closed by the coarser projection bound.
These experiments support a lower bound at both fractions, but they are not
certificates: the cuts and transcendental constants use floating arithmetic,
and local nonlinear optimization does not prove a global box maximum.

We independently replayed one local box at output fraction $0.22$.  The
checker uses rational dual potentials, directed bounds for every exponential
and logarithm, and a rational dual for a 65-variable linear program.  It also
checks the exact domain bounds and a rational composition in the box.  The
result certifies upper $-8.72\cdot10^{-6}$ on a coordinatewise
radius-$1/2000$ neighborhood of the numerical maximizer.  This certificate
does not cover the complete composition domain.

We also tested whether this local method scales to the full domain.  A global
outer approximation uses all degree-two composition moments, simplex moment
identities, first-level RLT inequalities, and separated covariance cuts.  Its
current partition has 10,948 cells after 11,414 splits.  The largest floating
upper bound is still $7.73\cdot10^{-3}$, although the fixed dual bank evaluates
to $-6.41\cdot10^{-3}$ at the reported composition.  The positive value is
therefore a relaxation gap: the lifted moments can mix products that do not
arise from one composition.

An explicit vertex-product LP reproduces this gap on the worst cell.  Its
optimum mixes seven vertex triples, even though the fixed dual bank is negative
at every triple.  This experiment suggests that additional convex moment
inequalities will not suffice.  A successful certificate must instead retain
which dual cut witnesses each part of the composition domain.  We treat this
conclusion as a diagnostic, not as a theorem about all possible relaxations.

We also tested a cut-selection relaxation that is stronger than a fixed
mixture.  On each chamber, affine nonnegative weights form a partition of
unity over cut pairs.  The resulting weighted sum of exponentiated cut bounds
is a valid pointwise upper bound.  Six coordinate comparisons partition the
worst cell into seven chambers.  Multistart continuous optimization gives a
negative upper bound on five chambers.  Two remain open.  For the harder open
chamber, an exchange procedure adds the best cut pair at each interior
violation; after 20 rounds the sampled constraints have margin
$-3.96\cdot10^{-4}$, while the largest violation found is
$2.64\cdot10^{-4}$.  This experiment localizes the remaining obstruction, but
does not certify a global maximum or the floating cut constants.

The partition itself is independently auditable.  A rational replay
reconstructs a 250,553-node prefix through split 10,100 and checks the next
1,314 recorded splits.  A separate 138-variable replay certifies upper
$-1.85\cdot10^{-2}$ on one radius-$0.01$ leaf.  These checks validate the
partition and the local arithmetic, but they do not prove the global
$0.22$ bound.  Closing that bound requires either further subdivision or a
stronger description of the rank-one product set.

\begin{conjecture}[First-moment threshold]\label{conj:first-moment}
For each fixed prime $p$, the minimum distance of the repetition-three BAA
ensemble satisfies
\[
  \frac{d_{\F_p}(E_{P_1,P_2})}{N}
  \xrightarrow{\;\Pr\;}\delta_{\rm FM}(p)
  \qquad\text{as }K\to\infty.
\]
\end{conjecture}

The first moment addresses only the lower side of
Conjecture~\ref{conj:first-moment}.  A negative bulk exponent, combined with
the spacing argument, excludes smaller codewords with probability tending to
one.  A positive exponent proves only that their expected number is
exponential.  Rare permutation pairs could dominate that expectation.  The
reverse inequality therefore requires concentration, a second-moment
argument, or another typical-existence theorem.  In particular, the numerical
quinary transition near $0.260$ is a target for certification, not a claimed
distance theorem.

There is, however, initial evidence that a second-moment argument may be
possible.  Count one-dimensional message subspaces rather than nonzero
messages; this removes the four scalar representatives of each quinary
codeword support.  A pair of distinct subspaces is described coordinatewise
over the additive group $\F_5^2$.  Its two accumulator terms are stationary
information payments over that group.  The following identity rules out one
potential source of pairwise correlation.

For a finite abelian group $G$, write $I_{\rm stat}^G(\alpha,\beta)$ for the
minimum of $I(D;S)$ over laws satisfying
$D\sim\alpha$, $S\sim\beta$, and $D+S\sim\beta$.

\begin{proposition}[Product tensorization]\label{prop:pair-tensor}
For distributions $\alpha$ and $\beta$ on $G$,
\[
 I_{\rm stat}^{G^2}(\alpha\otimes\alpha,\beta\otimes\beta)
 =2I_{\rm stat}^G(\alpha,\beta).
\]
\end{proposition}

\begin{proof}
The product of two minimizing one-coordinate laws proves the upper bound.
For the reverse bound, let $Q$ be any feasible law for the left-hand side
and group its variables as $X_i=(D_i,S_i)$.  If
$R_i=\alpha\otimes\beta$, the relative-entropy chain rule gives
\[
 D_{\rm KL}(Q_{X_1X_2}\,\|\,R_1\otimes R_2)
 =D_{\rm KL}(Q_{X_1}\,\|\,R_1)
  +D_{\rm KL}(Q_{X_2}\,\|\,R_2)
  +I_Q(X_1;X_2).
\]
Each coordinate projection of $Q$ is feasible for the one-coordinate
problem.  The first two terms are therefore at least
$I_{\rm stat}^G(\alpha,\beta)$, and the last term is nonnegative.
\end{proof}

Proposition~\ref{prop:pair-tensor} shows that two replicas cannot lower an
accumulator's information payment while their increment and state
compositions remain products.  A failure of the second moment must instead
come from nonproduct pair compositions.  We tested this possibility locally
at the quinary optimizer for output fraction $0.261$.  Holding both
one-coordinate marginals fixed leaves $16$ correlation directions in each of
the message, middle, and output compositions.  The resulting
$48$-dimensional numerical Hessian is negative definite: its eigenvalues
range from $-1.46\cdot10^{-2}$ to $-2.50\cdot10^{-5}$.  This computation is
not a certificate, but it finds no local correlation mode that improves the
pair exponent.  A proof of the conjecture still has to exclude separated
correlated maxima and control finite-type prefactors.

The spacing argument also gives a theorem for every fixed prime.

\begin{theorem}[Prime ladder]\label{thm:prime-ladder}
Fix a prime $p$ and an integer $m\ge4$.  If
\[
  288(p-1)^2m^4<(m-2)^6(m+2)^3,
\]
then
\[
  \Pr[d_{\F_p}(E_{P_1,P_2})\le\lfloor N/m^2\rfloor]
  \longrightarrow0.
\]
\end{theorem}

For every fixed $p$, a sufficiently large $m$ satisfies the inequality.
Hence this addition-only ensemble is asymptotically good over every fixed
prime field.  The resulting bound is of order $p^{-4/5}$ and becomes useless
when $p$ grows exponentially in the security parameter.

\subsection{What large characteristic proves}

The cryptographic regime instead has $N\approx2^{20}$ and
$\log_2p\approx128$.  Here $p$ is much larger than $N$, but its bit length is
much smaller than $N$.  Two deterministic statements explain why large
characteristic suppresses the observed local failures.

\begin{proposition}[Absolute distance floor]\label{prop:floor}
For every realized repetition-three BAA encoder,
\[
  d_{\F_p}\ge
  \min\left\{d_{\F_2},
    \left\lceil\log_{\sqrt{20}}p\right\rceil\right\}.
\]
\end{proposition}

The proof uses a sparse integer parity-check matrix.  Each row has Euclidean
norm at most $\sqrt{20}$.  If a support of size $d<d_{\F_2}$ becomes dependent
modulo $p$, an odd nonzero $d$-by-$d$ minor must be divisible by $p$.
Hadamard's inequality gives $p\le20^{d/2}$.  A prime near $2^{128}$ therefore
gives an unconditional distance floor of about $60$, not a linear distance.

\begin{proposition}[Cycle-rank barrier]\label{prop:cycles}
Suppose a connected residual component of the sparse encoder circuit has full
rank over $\mathbb Q$ but loses rank over $\F_p$.  Its cycle rank is at least
$\lceil\log_2p\rceil$.
\end{proposition}

To see the bound, expand a nonzero minor over perfect matchings.  After fixing
a spanning forest, the nonforest edges identify each matching.  A component
of cycle rank $\beta$ therefore contributes at most $2^\beta$ signed terms.
Divisibility by $p$ implies $p\le2^\beta$.

Propositions~\ref{prop:floor} and~\ref{prop:cycles} remove all known
small-characteristic traps.  They do not control an adversarial support that
creates a large, nearly square residual core.  The missing theorem is a
uniform random-saturation statement over every support of size below the
binary threshold.  Fixed-support experiments show substantial rank surplus,
but a minimum-distance witness selects its support after seeing both
interleavers.

\section{The Walsh--Hadamard alternative}\label{sec:wht}

Walsh--Hadamard PCGs take a different route.  They define a quasi-abelian code
directly over a large prime field and use a transform containing only additions
and subtractions~\cite{wht26}.  This avoids an unsupported binary-distance
transfer.  Random spectral multipliers still supply a large family of
projective types.

Removing those multipliers would make the encoder still cheaper.  The natural
sign-only variant, however, has a structural obstruction.  Let
$V=\F_2^r$, $N=2^r$, and let $c\ge2$ be the block count.  At each frequency
$\xi\in V$, write $\tau(\xi)\in\mathbb P^{c-1}(F)$ for the projective type of
the multiplier vector.

\begin{proposition}[Affine-flat obstruction]\label{prop:flat}
If $\tau$ is constant on an affine flat of dimension $s$, then the relative
minimum distance of the length-$cN$ code is at most $2^{-s}$.
\end{proposition}

\begin{proof}
On the flat, write the multiplier vector as a fixed projective vector times a
nonzero scalar.  Choose the message spectrum to cancel that scalar and equal a
character on the flat.  Its inverse Walsh transform is supported on one coset
of the orthogonal subspace.  Each nonzero output block has weight $N/2^s$.
\end{proof}

Suppose the types are independent and supported on at most $M$ values.  A
partition into parallel $s$-flats gives
\[
  \Pr[\Delta>2^{-s}]
  \le
  \exp\left(-\frac{N}{2^s}M^{1-2^s}\right).
\]
Choosing $s\asymp\log_2(\log N/\log M)$ yields
\[
  \Delta=O\left(\frac{\log M}{\log N}\right)
\]
with probability tending to one.  Every fixed spectral alphabet therefore
fails asymptotically.

This obstruction leaves a finite design space.  At $N=2^{20}$ and rate
$1/4$, an alphabet of $2^{10}$ signed small integers makes the simplest
dangerous three-flat union bound smaller than $2^{-128}$.  Multiplication by a
small integer can be replaced by an addition chain.  The calculation excludes
only one witness family; a global distance theorem for the growing-alphabet
variant remains open.

\section{Discussion and research agenda}\label{sec:discussion}

The investigation gives a clean division of the design space.  Dense binary
matrices prove that field-universal distance is compatible with positive rate.
BAA proves that near-linear addition-only encoding is compatible with positive
distance over every fixed prime.  Walsh--Hadamard codes prove that direct
large-field analysis is compatible with a multiplication-free transform.  No
current result combines all three properties at cryptographic parameters.

A PCG that deploys a binary or signed encoder over another field should make
one of the following justifications explicit.

\begin{enumerate}
  \item A scalar-extension argument applies because the target field extends
  the field used in the distance proof.
  \item A combinatorial certificate, such as unique-neighbor expansion,
  proves distance independently of the field.
  \item A direct analysis proves distance over the deployed field.
  \item The security statement tolerates a quantified distance loss, and its
  noise parameters are adjusted accordingly.
  \item The distance claim is labeled as a heuristic and supported by
  reproducible experiments rather than presented as a theorem.
\end{enumerate}

The most direct open problem is the following.

\begin{quote}
Construct a distribution of binary or signed $k$-by-$n$ matrices with
$k=\Theta(n)$, encoding complexity $O(n\log^c n)$, and a constant $\delta>0$
such that one sampled matrix has distance at least $\delta n$ over $\F_2$ and
over every prescribed $\lambda$-bit prime field, except with probability
$\negl(\lambda)$.
\end{quote}

For BAA, the missing step is uniform rank anti-concentration for the residual
cores selected by low-weight supports.  For Walsh--Hadamard codes, it is an
inverse theorem that classifies every low-weight spectrum under a growing
multiplier alphabet.  An intermediate construction may be easier: small dense
field-universal blocks could be composed through an expander or another sparse
mixing network.  Such a construction would trade some of BAA's speed for a
more transparent distance certificate.

The central conclusion is therefore positive but qualified.  Changing
characteristic is not harmless, and binary distance alone is not evidence.
Nevertheless, good addition-only codes over large prime fields are not ruled
out.  The remaining task is to make their mixing both fast and certifiable.

\bibliographystyle{splncs04}
\bibliography{references}

\end{document}
